% Abhakseite „Das kann ich“ – Zone nur im Gesamt (\mitzone)
%% TODO Ich-kann-Sätze: Platzhalter wie in den Aufgabendateien
\begin{abhakseite}
\ifdefined\mitzone
\abhakgruppe{Kennst du schon}
\abhak{1}{Ableitungen bilden}
\abhak{2}{Gleichungen lösen}
\abhak{3}{Charakteristische Punkte (Achsenschnitt-, Hoch-, Tief-, Wendepunkte) am Graphen ganzrationaler Funktionen ablesen und im Sachzusammenhang deuten}
\abhak{4}{Funktionswerte berechnen und Punktprobe}
\abhak{5}{Verlauf quadratischer und kubischer Graphen grob skizzieren (Öffnung, Symmetrie, Grenzverhalten)}
\abhak{6}{Ableitung als Tangentensteigung und lokale Änderungsrate deuten (Grundvorstellung)}
\abhak{7}{Gleichungen lösen – Fehler finden}
\abhak{8}{Gleichungen lösen – selbst rechnen}
\fi
\abhakgruppe{Einheit 1 · Monotonie und erste Ableitung}
\abhak{9}{Monotonie}
\abhak{10}{Monotonie – Fehler finden, Begründen, Anwendung}
\abhakgruppe{Einheit 2 · Extrempunkte}
\abhak{11}{Extrempunkte berechnen}
\abhak{12}{Extrempunkte nachweisen}
\abhak{13}{Abstand zweier Extrempunkte verschiedener Graphen mit einer Schranke vergleichen}
\abhak{14}{Abstand zweier Extrempunkte über die Punktsymmetrie berechnen}
\abhak{15}{Extrempunkte berechnen und ihre Verbindungsgerade als Winkelhalbierende nachweisen}
\abhak{16}{Nullstellen und Extremstelle einer Parabel berechnen}
\abhak{17}{Berührung des Graphen mit der x-Achse über die Extrempunkte begründen}
\abhak{18}{Größten Funktionswert über Monotonie, Grenzverhalten und Symmetrie begründen und Wertemenge angeben}
\abhak{19}{Extremstelle zwischen zwei Stellen mit gleichem Funktionswert ohne Rechnung begründen}
\abhak{20}{Lage eines Graphen zwischen zwei Geraden auf einem Intervall über die Differenzfunktionen nachweisen}
\abhak{21}{Aussagen über Funktions- und Ableitungswert nahe dem Tiefpunkt ohne Rechnung beurteilen}
\abhak{22}{Extrempunkte berechnen – Fehler finden, Begründen, Anwendung}
\abhakgruppe{Einheit 3 · Krümmung und Wendepunkte}
\abhak{23}{Wendepunkte und Krümmung}
\abhak{24}{Gerade durch die beiden Wendepunkte aufstellen und parallele Gerade mit genau einem gemeinsamen Punkt einzeichnen}
\abhak{25}{Punktprobe am Graphen}
\abhak{26}{Wendestelle nachweisen und Winkel der Wendetangente mit der x-Achse über die Steigung −1 zeigen}
\abhak{27}{Wendestellen einer Sinusfunktion als ganzzahlig nachweisen und die beiden Wendetangentensteigungen zeigen}
\abhak{28}{Anzahl der Schnittpunkte von Geraden durch den Wendepunkt mit dem Graphen nach der Steigung unterscheiden}
\abhak{29}{Wendepunkte und Krümmung – Fehler finden, Begründen, Anwendung}
\abhakgruppe{Einheit 4 · Graph und Ableitungsgraph}
\abhak{30}{Graph und Ableitungsgraph}
\abhak{31}{Gemeinsamen Punkt zweier Graphen über gleiche Flächeninhalte indirekt begründen}
\abhak{32}{Logarithmus einer Exponentialfunktion als lineare Funktion nachweisen und Steigung und Achsenabschnitt angeben}
\abhak{33}{Graph und Ableitungsgraph – Fehler finden, Begründen, Darstellungswechsel}
\abhakgruppe{Einheit 5 · Kurvenuntersuchung im Sachzusammenhang}
\abhak{34}{Sachzusammenhang}
\abhak{35}{Funktionswert im Sachzusammenhang deuten}
\abhak{36}{Mittleren Wert einer periodisch schwankenden Größe am Graphen ablesen}
\abhak{37}{Sachzusammenhang – Fehler finden, Begründen, Anwendung, Darstellungswechsel}
\end{abhakseite}
